

\chapter{Submanifolds and Hypersurfaces}
\label{ch:hypersurface}

\section{Introduction}

This is a crucial ingredient in understanding almost anything that has to do with Numerical Relativity, the art of solving EFEs by the help of numerical methods. I will closely follow the treatment given by David Malament in his book on the foundations of GR and Newtonian gravity. In my opinion, among all the GR textbooks, this one has the most systematic treatment of this topic. Though, being written by a philosopher, the main idea of the book is not to teach differential geometry so that one can use it in research, rather it is to present only that much which forms the core of our understanding of the topic. However, this much is actually sufficient for diving deep into a proper numerical relativity textbook.



\section{Imbedding}

 Let S and M be two manifolds. Let there be a smooth map $\Psi:S \longrightarrow M $. $\Psi$ is called an imbedding if the following three conditions hold.
 \begin{itemize}
     \item $\Psi$ is injective.
     \item at all p in S, the push forward map $\lera{\Psi_p}_*: T_p S \longrightarrow T_pM$ is injective.
     \item the inverse map $\Psi^{-1}: \Psi \left[ S \right] \longrightarrow S$ is continuous with respect to the relative topology in $\Psi\left[S \right]$.
 \end{itemize}


At first glance, this definition might seem unmotivated and built out of thin air. \textcolor{red}{TODO: learn and write the reasons behind this definition}

Now, suppose that S$\subseteq$M, then we say that S is an imbedded submanifold of M if the map $id:S \longrightarrow M $ is an imbedding. If the dimension of M is n, and the dimension of S is k, and $k=n-1$ then we say that S is a \textit{hypersurface} in M. 


In this and the next few sections, let me fix S to be a k-dimensional imbedded submanifold of M. Later on we will think more specifically in terms of hypersurfaces. But for now let me go through some mathematical machinary needed for imbedded submanifolds in general. 



\section{S-tensors vs M-tensors}

At a point p, we can talk about tensors with respect to the atlas of M, or with respect to the atlas of S. We refer to the former as M-tensors, and the later as S-tensors. The S-tensors are S's own objects and we do not need to refer to M at all while talking about them. We will use a tilde over the name of the tensor to show that it is an S-tensor. 

Given S to be a metric submanifold of M ( I will explain what a metric submanifold is in a bit), we can divide $T_pM$ into two subspaces, one would be a subspace of vectors ``tangent'' to S, and one ``normal'' to S. 

So what, then, is an M-vector at p that is ``tangent'' to S? (I am using quotes because this tangent is not related to the word ``tangent'' in ``tangent vector''). 

\begin{definition}
    A vector $X\in T_pM$ is \textit{tangent} to S if $X\in (id_p)_*(T_p S)$.
\end{definition}


\begin{definition}
    A covector $\eta \in T^*_pM$ is normal to S if for all vectors $X \in T_pM $ that are tangent to S, $\eta(X)=0$
\end{definition}

We denote the space of all M-vectors tangent to S at p as $\perp_p$ and the space of all M-covectors normal to S at p as $N^*_p$. As we took $T_pS$ and then pushed it forward with an imbedding (which basically one to one maps $T_pS$ to $T_pM$, the dimension of $\perp_p$ is k. We can also show that the dimension of $N^*_p$ is $n-k$.

Another way to describe tangent vector comes from what it means to push forward a vector from one manifold to another. For a vector $ \tilde{X}$ in S, there is obviously a curve $ \gamma$ to which the vector $ \tilde{X}$ is tangent to. Since push forward of a vector tangent to a curve is tangent to the image of the curve, $ id_* \tilde{X}$ is tangent to $ id( \gamma)$. In fact, we are considering S to be a subset of M, and $ id( \gamma)$ is actually just $ \gamma$. This means, any M-vector at p that is tangent to a curve $ \gamma$ which lies entirely in S, is a memeber of $ \perp _p$.
Wait, how do you show the converse.....?????


\begin{proposition}
    The dimension of $N^*_p$ is $n-k$.
\end{proposition}

\begin{proof}
    First of all, we know that dimension of $\perp_p$ is k. Take k vectors $\overset{1}{X}, \overset{2}{X},....,\overset{k}{X}$ in $\perp_p$ that are linearly independent. Now, to make a basis of $T_pM$ we need some other linearly independent vectors outside of $\perp_p$. Let's say the total basis is   $\overset{1}{X}, \overset{2}{X},....,\overset{k}{X},\overset{k+1}{X},....,\overset{n}{X}$. Now, let us take the dual basis of the previous basis : $\overset{1}{\eta}, \overset{2}{\eta},....,\overset{k}{\eta},\overset{k+1}{\eta},....,\overset{n}{\eta}$. Then, the covectors $\overset{k+1}{\eta},....,\overset{n}{\eta}$ span a space $N^{'}$ with the following property: any member $\eta$ of $N^{'}$, being a combination of $\overset{k+1}{\eta},....,\overset{n}{\eta}$ , will actually produce $\eta(X)=0$ for any member X of $\perp_p$ as X is a combination of $\overset{1}{X}, \overset{2}{X},....,\overset{k}{X}$. That means the space $N^{'}$. Clearly, $N^{'}$ is actually what we are calling $N^*_p$. And the dimension is $n-l$
\end{proof}

Once we have $N^*_p$ defined, we can prove the following proposition, which will be very useful later on where we use this result at the same rate, or even more, than the first definition of a vector being tangent. 


\begin{proposition}
    An M-tensor X is tangent to S at p iff for any $\eta \in N^*_p$ , $\eta(X)=0$.
\end{proposition}

\begin{proof}
    Let us use the basis vectors as defined before, that is, 
    $T_pM$ is spanned by $\{ \overset{1}{X}, \overset{2}{X},....,\overset{k}{X},\overset{k+1}{X},....,\overset{n}{X}\}$ where the first k spans $\perp_p$. Let $T^*_pM$ be spanned by the dual vectors $\{\overset{1}{\eta}, \overset{2}{\eta},....,\overset{k}{\eta},\overset{k+1}{\eta},....,\overset{n}{\eta}\}$ where the last $n-k$ vectors span $N^*_p$.

    Now, let $X \in \perp_p$. \\
    This means X has the form \[X= \sum_{1}^{k} a_j \overset{j}{X}\] For any $\eta \in N^*_p$, $\eta$ has the form, \[\eta = \sum_{k+1}^{n} b_j \overset{j}{\eta}\] Clearly, $\eta(X)=0$.


    Lets now assume that we have an $X \in T_pM$ such that for any $\eta \in N^*_p$, $\eta(X)=0$.\\
    This means in the expansion of X, there cannot be any term of the form $a_j \overset{j}{X}$ where $j>k$. And so X is a linear combination of only the basis vectors $\{\overset{1}{X},.... \overset{k}{X}\}$. And that means, $X \in \perp_p$
    
    
\end{proof}

While the M-vectors tangent to S come directly from S via a push forward map, the M-covectors pulled back to S actually give zero. 

\begin{proposition}
     $\eta \in N^*_p \iff (id_p)^*(\eta)=0$
\end{proposition}
\begin{proof}
    Let $\eta \in N^*_p$ \\
    Let $\Tilde{X} \in T_pS $. So, $\left( (id_p)^*(\eta) \right)\lera{\Tilde{X}} = \eta \lera{\lera{id_p}_*\Tilde{X}}=0$. Since this is true for any $\Tilde{X}$ in $T_pS$, we can conclude that $(id_p)^*(\eta)=0$.

    Now, let $(id_p)^*(\eta)=0$ where $\eta \in T_p^*M $\\
    For any $\Tilde{X} \in T_pS$,
    \begin{align*}
        &\lera{(id_p)^*(\eta)}(\Tilde{X}) =0 \\
        \implies & \eta \lera{(id_p)_*(\Tilde{X})} =0 \\
    \end{align*}
    this means $\eta$ acting on any member of $\perp_p$ is zero. And that means $\eta \in N^*_p$.
    
\end{proof}


Now, more general M-tensors at p can also be tangent and normal, in different indices. For example, we can say when a (2,3) M-tensor $T^{ab}_{cde}$ at p is tangent in the first upper index a, or normal in the third lower index e. The definition that will follow can be easily generalized to the most general case, but I will only give a restricted special case definition, because writing down the most general one is more annoying than illuminating.
\begin{definition}
    A (2,3) M-tensor $\tensor{T}{ab}{\quad cde}$ at point p is tangent to S in the second upper index b if for any $\eta \in N^*_p$, $\tensor{T}{ab}{\quad cde} \eta_b =0$. It is normal to S in the third lower index e if for any $X \in \perp_p$ , $\tensor{T}{ab}{\quad cdn}X^n=0$. And so on for tensors of different index structures.
\end{definition}

So far we are only considering being tangent on an upper index, and being normal on a lower index. But it is also possible to be normal on an upper index and to be tangent on a lower index if the manifold M is a metric manifold (even at this point I do not need the submanifold to have a metric).

\begin{definition}
    $\tensor{T}{ab}{\quad cde}$ is normal in the first upper index a if $ T_{a\;\;\;cde}^{\;\;b} = \tt{g}{}{am}\tt{T}{mb}{cde}$ is normal in the index a. Similarly T is tangent in the second lower index d if $ T^{ab\;\;d}_{\;\;\;\;c\;\;\;e}=g^{dm}\tensor{T}{ab}{\quad cme}$ is tangent in the index d. And so on for tensors of other index structure. 
\end{definition}

Basically, this definition says, we need to raise or lower an index. Then if after raising a lower index the tensor is tangent, then it is tangent in that index in it's lower form too. And if after lowering an index, the tensor is normal, then that tensor is normal in that index in the raised form too.

Right now, for M-vectors and M-covectors at p, we have 4 different spaces in terms of tangency and normalcy to S. 

\begin{itemize}
    \item $\perp_p$ is the space of M-vectors tangent to S at p.
    \item $\perp^*_p$ is the space of M-covectors tangent to S at p.
    \item $N_p$ is the space of M-vectors normal to S at p.
    \item $N^*_p$ is the space of M-covectors normal to S at p.
\end{itemize}


We know that the dimension of $\perp_p$ is k and the dimension of $N^*_p$ is $(n-k)$. What is the dimension of $\perp^*_p$ ? Or $N_p$ ?

\begin{proposition}
    The dimesion of $\perp^*_p$ is k and the dimension of $N_p$ is $(n-k)$.
\end{proposition}

\begin{proof}
    We know that the metric tensor $g_{ab}$ creates a bijective mapping between $T_pM$ and $T^*_pM$. So thinking from that perspective, $\perp^*_p$ is basically the set of all covectors $X$ such that $X_a = g_{am} X^m$ where $X \in \perp_p$. And the dimension of such a space is obviously the dimension of $\perp_p$, that is, k.

    In a similar way, the space $N_p$ is made of all vectors $\eta$ that has the form $\eta^a = g^{am} \eta_m$ where $\eta \in N^*_p$. And the dimension of such a space is equal to the dimension of $N^*_p$, which is $(n-k)$.
\end{proof}




Now we might be thinking of a nice possibility. Maybe, $T_pM = \perp_p + N_p $















\section{Sticking to Only M-Tensors}

At any point p, there is a natural isomorphism between S-tensors of any index structure and M-tensors of the same index structure that are tangent to S. This is the very basis of how we talk about a hypersurface. The idea is that we will commit to M-tenosrs tangent to S and not talk about S-tensors.

Let me first describe the isomorphism. Here I will satisfy myself by only defining the map, and I will skip the proof of the map being an isomorphism. 

The map $ \phi_{p}$ takes an S-tensor $ \tt{ \tilde{X}}{ab}{cd}$ to $ \phi_{p} \lera{ \tt{ \tilde{X}}{ab}{cd}}$, which is an M-tensor. We define the isomorphism in stages, first for scalars, then for vectors, covectors and finally for tensors of arbitray index structure. 

First, for a scalar $ \alpha$ at p, $ \phi_{p} \lera{ \alpha}= \alpha$. There is no need to use tilde as scalars are just numbers. 

Second, for an S-vector $ \tt{ \tilde{X}}{a}{}$ at p, $ \phi_{p} \lera{ \tt{ \tilde{X}}{a}{}} = \lera{ id_p}_* \tt{ \tilde{X}}{a}{}$. By the very definition of a vector being tangent to S at p, $ \phi_p \lera{ \tt{ \tilde{X}}{a}{}}$ is tangent to S. 

Third, for a S-covector $ \tt{ \tilde{ \eta}}{}{a} $ at p, $ \phi_{p} \lera{ \tt{ \tilde{ \eta}}{}{a}}$ is defined by its action on any M-vector $ X^a$. 
 
\[
\phi_{p} \lera{ \tt{ \tilde{ \eta}}{}{a}} \tt{X}{a}{}=
\begin{cases}
 \tt{ \tilde{ \eta}}{}{a} \left(\phi_{p}^{-1} X^a \right); &  X \in \perp_p \\
0; & X \in N_p \\
\end{cases}
\]


We remark that this is well defined, as $ \phi_{p}^{-1} \tt{X}{a}{}$ is well defined.  Also, notice that the second clause of the definition of $ \phi_p \lera{ \tt{\tilde{ \eta}}{}{a}}$ ensures that it is tangent to S.


Lastly, lets take an arbitray index structure S-tensor $ \tt{ \tilde{X}}{ab}{c}$, then,

\[
\phi_{p} \lera{ \tt{ \tilde{X}}{ab}{c}} \tt{A}{}{a} \tt{ B}{}{b} \tt{ C}{c}{} = 
\begin{cases}
    \tt{ \tilde{X}}{ab}{c} \lera{ \phi_{p}^{-1}A_a} \lera{ \phi_{p}^{-1} B_b} \lera{ \phi_{p}^{-1}C^c } ; & A_a, B_b ,C^c \text{ are tangent to S}\\
    0 ; & \text{if } A_a, B_b ,\text{or }C^c  \text{ is normal to} \\ & \text{S}\\
\end{cases}
\]

Again, $ \phi_p \lera{ \tt{ \tilde{X}}{ab}{c}}$ is tangent to S by the second clause above. 

One can also show that $ \phi_p$ commutes with all tensor operations, i.e. addition, outer multiplication, index substitution, contraction.

For example,

\begin{align*}
    \phi_p \lera{ \tt{ \tilde{X}}{a }{} \tt{ \tilde{Y}}{b}{}} = \phi_p \lera{ \tt{ \tilde{X}}{a}{}} \phi_{p} \lera{\tt{ \tilde{Y}}{b}{}} 
\end{align*}



In short, at any point p in S, the tensor algebra of S-tensors is isomorphic to the tensor algebra of M-tensors that are tangent to S. With every S-tensor, I get an evil twin M-tensor of the same index structure.

Now, this can be extended to the whole of S, and we have a map $ \phi$ that takes any S-tensor field to an M-tensor field that is tangent to S and have the same index structure, and they also follow the same tensor algebra on both sides of the ``river''.

\subsection*{S-smoothness and M-smoothness}

First of all, let us recall  a very handy notion of the smoothness of a tensor field. Say, we have a manifold M, and an atlas $ \mathcal{A}$. Now, a tensor field $ \tt{T}{a}{b}$ is smooth iff the components are smooth functions with respect to any coordinatization of the manifold with charts from the atlas $ \mathcal{A}$. This means the idea of smoothess is tied to the atlas of a manifold. 

Now think of the manifold S. Let me define by O an open set of M that covers S, i.e. S is a subset of O. This is going to be our working definition of O for quite some time in this discussion, unless otherwise specified. An M-tensor field T tangent to S can be smooth in two senses. If an extension of T on O, named $ \overset{+}{T}$, is smooth wrt $ \mathcal{A}_M$, then T is M-smooth. If, on the other hand, the S-tensor field $ \phi^{-1} T$ is smooth wrt $ \mathcal{A}_S$, then T is S-smooth. Now, there is a theorem that says for any M-tensor field T tangent to S, it is M-smooth iff it is S-smooth. So saying that T (when it is tangent to S) is smooth means it is smooth in both ways. On the other hand if  T is not tangent to S, then it can only be M-smooth.




\subsection*{ $ \tt{h}{}{ab}$}

This is the star of the show. We call it a projector for reasons to become clear soon. Let us define it first. Say there is a metric $ \tt{g}{}{ab}$ on M. We define the ``induced metric'' or ``first fundamental form'' on S by taking the pull back of $ \tt{g}{}{ab}$ wrt to the identity map. The induced metric, an S-field, $ \tt{ \tilde{h}}{}{ab} = id^* \tt{g}{}{ab}$. Then, we define the smooth symmetric M-tensor $ \tt{h}{}{ab}$ which is tangent to S as the $ \phi$-image of this induced metric $ \tt{\tilde{h}}{}{ab}$, so that $ \tt{ h}{}{ab }= \phi \lera{ \tt{ \tilde{h}}{}{ab}}= \phi \lera{ id^* \tt{g}{}{ab}}$. 

This is as counterintuitive and hard to remember as possible. So we actually use an equivalent definition of $ \tt{h}{}{ab}$. Here, we define $ \tt{h}{}{ab}$ by the way it acts on M-vectors. 

\begin{align*}
    \tt{h}{}{ab} \tt{X}{a}{} \tt{Y}{b}{} = 
    \begin{cases}
        \tt{g}{}{ab} \tt{X}{a}{} \tt{Y}{b}{} ; & \text{if } X,Y \text{ are tangent to S}\\
        0; & \text{if any of } X,Y \text{ are normal to S} 
    \end{cases} \numberthis \label{hab_defiend} \\
\end{align*}


Let me show that the equivalence works. Lets say both X and Y are tangent to S. Then,

\begin{align*}
    \tt{h}{}{ab} \tt{X}{a}{} \tt{Y}{b}{} &= \phi \lera{id^* \tt{g}{}{ab}} \tt{X}{a}{} \tt{Y}{b}{} \\
    &= \lera{ id^* \tt{g}{}{ab}} \lera{ \phi^{-1} \tt{X}{a}{}} \lera{ \phi^{-1} \tt{Y}{b}{}} \\
    &= \tt{g}{}{ab} \; id_* \lera{ \phi^{-1} \tt{X }{a}{}} \; id_* \lera{ \phi ^{-1} \tt{Y}{b}{}} \\
    &= \tt{g}{}{ab} \; \phi \lera{ \phi^{-1} \tt{X}{a}{}} \phi \lera{ \phi^{-1} \tt{Y}{b}{}} \\
    &= \tt{g}{}{ab} \tt{X}{a}{} \tt{Y}{b}{} \\
\end{align*}




If any of X or Y is not tangent, then by virtue of $ h_{ab}$ being tangent to S, $ h_{ab} X^a Y^b$ is 0.

We define a perpendicular projection tensor $ k_{ab}$ too which is a smooth symmetric M-tensor field.

\begin{align*}
    k_{ab} = g_{ab} - h_{ab} \numberthis \label{kab_defined}\\
\end{align*}


Notice that $ k_{ab}$ is normal to S. If X is a M-vector tangent to S, then, $ k_{ab} X^a = g_{ab} X^a - h_{ab} X^a = g_{ab} X^a - g_{ab} X^a =0 $. So $ k_{ab}$ is normal to S in the first index. But being a symmetric tensor, it must be also normal in the second index too, so it is normal as a whole. 



The following propositions show that it is okay to call $ h_{ab}$ a tangential projector and $ k_{ab}$ a perpendicular or normal projector.  notice that $ k_{ab}$ is normal to S. 





\begin{proposition}
     $ X^a$ is tangent to S $ \iff$ $ \tt{h}{a}{b} X^b = X^a \iff \tt{k}{a}{b} X^b =0 $
\end{proposition}



\begin{proof}
    We will show the first equivalence and then the second equivalence.
    
    First equivalence: 
    
    Let X be tangent to S. 

    I will show that the action of $X^a$ and $ \tt{h}{a}{b} X^b$ is the same on any M-covector $ \eta_a$. If $ \eta_a $ is normal to S, then both sides give zero. If $ \eta_a$ is tangent to S, then $ \tt{h}{a}{b} X^b \eta_a = h_{ab}  \eta^a X^b = g_{ab}  \eta^a X^b= X^a \eta_a$.


    Conversely, let $ \tt{h}{a}{b} X^b = X^a$.
    I need to show that $ X^a$ is tangent to S. Take $ \eta_a$ normal to S. Then $ \eta_a X^a = \eta_a \tt{h}{a}{b} X^b = 0 $. The last equality follows from the fact that $ h_{ab}$ is tangent in both indices.

    This completes the proof for the first equivalence. 
    
    Second equivalence: 

    Let $ \tt{h}{a}{b} X^b = X^a$.

   Then, $ \tt{k}{a}{b} X^b = \lera{ \tt{g}{a}{b}- \tt{h}{a}{b}} X^b = X^a- X^a =0$

    Conversely, let, $ \tt{k}{a}{b} X^b =0 $.

    Then, $ \tt{h}{a}{b } X^b = \lera{ \tt{g}{a}{b} - \tt{k}{a}{b}} X^b = X^a - 0 = X^a $

    This completes the proof of the second equivalence. 
    
    
    
\end{proof}



\begin{proposition}
   $ X^a$ is normal to S $ \iff$ $ \tt{k}{a}{b} X^b = X^a \iff \tt{h}{a}{b} X^b =0 $ 
\end{proposition}


\begin{proof}
    Similar as before. 
\end{proof}


While I stated proposition 4.9 and 4.10 for a (1,0) tensor field, i.e. a vector field $ X^a$, these propositions actually hold for tensor fields of arbitrary index structure. For example, if $ \tt{T}{ab}{c}$ is tangent to S in the second upper index, then $ \tt{h}{b}{m} \tt{T}{am}{c}= \tt{T}{ab}{c}$ and $ \tt{k}{b}{m} \tt{T}{am}{c}=0$.

\begin{proposition}
    $ \tt{h}{a}{b} \tt{h}{b}{c} = \tt{h}{a}{c}$
\end{proposition}

\textcolor{blue}{Let me pause here for a moment, and fill up the rest of the proofs up to the Gauss Codazzi equations}

\subsection*{Extrinsic Curvature}

We know that, 

\begin{align*}
    \tt{h}{m}{a} \tt{h}{n}{b} \tt{h}{p}{c} \cov{ \tt{h}{}{np}}{m} =0 \\
    \tt{h}{m}{a} \tt{k}{n}{b} \tt{k}{p}{c} \cov{ \tt{h}{}{np}}{m} =0 \\
\end{align*}


However, the hhk projection, $\tt{h}{m}{a} \tt{h}{n}{b} \tt{k}{p}{c} \cov{ \tt{h}{}{np}}{m} $ is not zero and we define a new quantity, called the extrinsic curvature of S, as follows,

\begin{align*}
    \tt{K}{}{abc} = \tt{h}{m}{a} \tt{h}{n}{b} \tt{k}{p}{c} \cov{ \tt{h}{}{np}}{m} \\
\end{align*}


$ \tt{K}{}{abc}$ has particular geometric significance, and we will show some of those significance in this section. 




\subsubsection*{Gauss Codazzi Equations}

Before proving those equations, let me remind one useful way of characterizing the Riemann tensor. If $ v^a$ is a vector field, then,

\begin{align*}
    [ \nabla_c \;, \nabla_d  ] v^a &= \tt{R}{a}{mcd} v^m  \\
    2 \nabla_{[c} \nabla_{d] } v^a &= \tt{R}{a}{mcd} v^m \\
    \nabla_{[c} \nabla_{d] } v^a &=  \frac{1}{2} \tt{R}{a}{mcd} v^m \\
\end{align*}


This can also be generalized, in the following way. Say, $ \tt{T}{a}{b}$ is a (1,1) field. Then the antisymmetrized double covariant derivative on it can be expanded in temrs of the Riemann tensor as follows, 

\begin{align*}
    2 \nabla_{[c} \nabla_{d]} \tt{T}{a}{b} &= \tt{R}{a}{mcd} \tt{T}{m}{b}- \tt{R }{m}{bcd} \tt{T}{a}{m} \\
\end{align*}

One final fact about the Riemann tensors that needs a reminder is, $ \tt{R}{}{abcd}= \tt{R}{}{cdab}$.



We will use a script R, $ \mathcal{R}$, to denote the Riemann tensor defined wrt to D. This is going to be a tangent tensor field. So, we define the Riemann tensor to be, $ \tt{ \mathcal{R}}{a}{mcd} v^m = 2 D_{[c} D_{d]} v^a $ where $ v^a $ is an M-field tangent to S.

\textcolor{blue}{To Do: Write down some more words about the proper definition of $ \mathcal{R} $}


Now, 
\begin{align*}
    D_{[c} D_{d]} v^a &= \tt{h}{m}{[c} \tt{h}{n}{d]} \tt{h}{a}{p} \cov{D_n v^p}{m} \\
    &= \tt{h}{m}{[c} \tt{h}{n}{d]} \tt{h}{a}{p} \nabla_m \lera{ \tt{h}{q}{n} \tt{h}{p}{r} \nabla_q v^r} \\
    &= \tt{h}{m}{[c} \tt{h}{n}{d]} \tt{h}{a}{p} \lera{ \nabla_m \tt{h}{q}{n}} \tt{h}{p}{r} \nabla_q v^r + \tt{h}{m}{[c} \tt{h}{n}{d]} \tt{h}{a}{p} \tt{h}{q}{n} \lera{ \nabla_m \tt{h}{p}{r}} \nabla_q v^r \\
    & + \tt{h}{m}{[c} \tt{h}{n}{d]} \tt{h}{a}{p} \tt{h}{q}{n} \tt{h}{p}{r} \nabla_m \nabla_q v^r \\
    &= \tt{h}{m}{[c} \tt{h}{n}{d]} \tt{h}{a}{r} \lera{ \nabla_m \tt{h}{q}{n} } \nabla_q v^r + \tt{h}{m}{[c} \tt{h}{q}{d]} \tt{h}{a}{p} \lera{ \nabla_m \tt{h}{p}{r}} \nabla_q v^r \\
    &+ \tt{h}{m}{[c} \tt{h}{q}{d]} \tt{h}{a}{r} \nabla_m \nabla_q v^r \\\
    & \lera{\text{ Mostly using  $ \tt{h}{a}{b} \tt{h}{b}{c}= \tt{h}{a}{c}$}} \numberthis \label{weird1}\\
\end{align*}



Now, let's look at them term by term. 
\newline First Term:

\begin{align*}
  & \tt{h}{m}{[c} \tt{h}{n}{d]} \tt{h}{a}{r} \lera{ \nabla_m \tt{h}{q}{n} } \nabla_q v^r \\
  &= \tt{h}{m}{[c} \tt{h}{n}{d]} \tt{h}{a}{r} \lera{ \nabla_m \tt{h}{}{nq} } \nabla^q v^r \\
  &= \tt{h}{a}{r} \tt{h}{m}{[c} \tt{h}{n}{d]}  \lera{ \nabla_m \tt{h}{}{nq} } \nabla^q v^r \\
  &= - \tt{h}{a}{r} \tt{K}{}{[cd]q} \nabla^q v^r \\
  & \lera{ \text{Using the relation $ \tt{K}{}{abc}= - \tt{h}{m}{a } \tt{h}{n}{b} \nabla_m h_{nc}$}}\\
  &= 0 \\
  & \lera{ \text{ Since the antisymmetrization of the first two indices of $ \tt{K}{}{abc}$ gives zero.}}
\end{align*}

 Second Term:

\begin{align*}
    & \tt{h}{m}{[c} \tt{h}{q}{d]} \tt{h}{a}{p} \lera{ \nabla_m \tt{h}{p}{r}} \nabla_q v^r \\
    &= \tt{h}{m}{[c} \tt{h}{q}{d]} \tt{h}{ap}{} \lera{ \nabla_m \tt{h}{}{pr}} \nabla_q v^r \\
    &= \tt{h}{m}{[c} \tt{h}{q}{d]} \tt{h}{pa}{} \lera{ \nabla_m \tt{h}{}{pr}} \nabla_q v^r \\
    &= \tt{g}{sa}{}\tt{h}{m}{[c} \tt{h}{q}{d]} \tt{h}{p}{s} \lera{ \nabla_m \tt{h}{}{pr}} \nabla_q v^r \\
    &= \tt{g}{sa}{}\tt{h}{m}{[c} \tt{h}{p}{|s|} \tt{h}{q}{d]}  \lera{ \nabla_m \tt{h}{}{pr}} \nabla_q v^r \\
    &= -\tt{g}{sa}{} \tt{K}{}{[c|sr|} \tt{h}{q}{d]} \nabla_q v^r \\
    &= -\tt{g}{sa}{} \tt{K}{}{[c|sr|} \tt{h}{q}{d]} \nabla_q \lera{ \tt{h}{r}{t} v^t} \\
    & \lera{ \text{The trick is to realize that v is tangent to S.}} \\
    &= - K^{\;\;a}_{[c \;\;\; |r|} \tt{h}{q}{d]} \nabla_q \lera{ \tt{h}{r}{t} v^t} \\
    &= - K^{\;\;a}_{[c \;\;\; |r|} \tt{h}{q}{d]} \lera{ \nabla_q \tt{h}{r}{t}} v^t - K^{\;\;a}_{[c \;\;\; |r|} \tt{h}{q}{d]} \tt{h}{r}{t} \nabla_q v^t \\
    &= - K^{\;\;a}_{[c \;\;\; |r|} \tt{h}{q}{d]} \lera{ \nabla_q \tt{h}{r}{t}} v^t -0 \\
    & \lera{
    \begin{aligned}[]
        &\text{The second term goes to zero as the extrinsic curvature $ \tt{K}{}{abc}$ is} \\
        &\text{normal in the third index.}
    \end{aligned}
    }\\
    &=- K^{\;\;a}_{[c \;\;\; |r|} \tt{h}{q}{d]} \lera{ \nabla_q \tt{h}{r}{t}} \lera{ \tt{h}{t}{u} v^u} \\
    &= - K^{\;\;ar}_{[c} \tt{h}{q}{d]} \lera{ \nabla_q \tt{h}{}{tr}} \lera{ \tt{h}{t}{u} v^u} \\
    &= -K^{\;\; ar}_{[c} \tt{h}{q}{d]} \tt{h}{t}{u} \lera{ \nabla_q h_{tr}} v^u \\
    &= K^{\; \; ar}_{[c} K_{d]ur} v^u \\
\end{align*}


Third Term:

\begin{align*}
    & \tt{h}{m}{[c} \tt{h}{q}{d]} \tt{h}{a}{r} \nabla_m \nabla_q v^r \\
    &= \tt{h}{m}{c} \tt{h}{q}{d} \tt{h}{a}{r} \nabla_{[m} \nabla_{q]} v^r \\
    &= \tt{h}{m}{c} \tt{h}{q}{d} \tt{h}{a}{r} \left(\frac{1}{2} \right) \tt{R}{r}{nmq} v^n \\
    &= \frac{1}{2} \tt{R}{r}{nmq} \tt{h}{m}{c } \tt{h}{q}{d} \tt{h}{a}{r} v^n \\
\end{align*}

Collecting the terms, I get from \cref{ weird1 },

\begin{align*}
    D_{[c} D_{d]}v^a &= \frac{1}{2} \tt{R}{r}{nmq} \tt{h}{m}{c } \tt{h}{q}{d} \tt{h}{a}{r} v^n +K^{\; \; ar}_{[c} K_{d]ur} v^u \\
    &= \frac{1}{2} \tt{R}{r}{mnq} \tt{h}{n}{c } \tt{h}{q}{d} \tt{h}{a}{r} v^m +K^{\; \; ar}_{[c} K_{d]mr} v^m \\
\end{align*}

Now, from our earlier characterization of $ \tt{ \mathcal{R}}{}{} $, we have,

\begin{align*}
    &\tt{ \mathcal{R}}{a}{mcd} v^m = 2 D_{[c} D_{d]}v^a \\
    \implies&  \tt{ \mathcal{R}}{a}{mcd} v^m =  \tt{R}{r}{mnq} \tt{h}{n}{c } \tt{h}{q}{d} \tt{h}{a}{r} v^m +2K^{\; \; ar}_{[c} K_{d]mr} v^m \\
\end{align*}


Now, taking any vector $ \eta$, not necessarily tangent, and then taking the tangential projection of that vector, $ v^m= \tt{h}{m}{s} \eta^s$,

\begin{align*}
    &\tt{ \mathcal{R}}{a}{mcd} \tt{h}{m}{s} \eta^s =\tt{R}{r}{mnq} \tt{h}{n}{c } \tt{h}{q}{d} \tt{h}{a}{r} \tt{h}{m}{s} \eta^s +2K^{\; \; ar}_{[c} K_{d]mr} \tt{h}{m}{s} \eta^s\\
    & \tt{ \mathcal{R}}{a}{mcd} \tt{h}{m}{s} \eta^s = \left( \tt{R}{r}{mnq} \tt{h}{n}{c } \tt{h}{q}{d} \tt{h}{a}{r} \tt{h}{m}{s}  +2K^{\; \; ar}_{[c} K_{d]mr} \tt{h}{m}{s}\right) \eta^s\\
    &\tt{ \mathcal{R}}{a}{mcd} \tt{h}{m}{s}  =  \tt{R}{r}{mnq} \tt{h}{n}{c } \tt{h}{q}{d} \tt{h}{a}{r} \tt{h}{m}{s}  +2K^{\; \; ar}_{[c} K_{d]mr} \tt{h}{m}{s}\\
    & \tt{ \mathcal{R}}{a}{scd}  =  \tt{R}{r}{mnq} \tt{h}{n}{c } \tt{h}{q}{d} \tt{h}{a}{r} \tt{h}{m}{s}  +2K^{\; \; ar}_{[c} K_{d]sr} \\
    & \lera{ \text{We assumed that $ \tt{ \mathcal{R}}{}{}$ is a tangent field, also K is tangent in the second index. }} \\
    & \tt{ \mathcal{R}}{a}{bcd}  =  \tt{R}{r}{mnq} \tt{h}{n}{c } \tt{h}{q}{d} \tt{h}{a}{r} \tt{h}{m}{b}  +2K^{\; \; ar}_{[c} K_{d]br} \\
    & \lera{ \text{$ s \rightarrow b $}} \\
    & \tt{ \mathcal{R}}{a}{bcd}  =  \tt{R}{m}{rnq} \tt{h}{n}{c } \tt{h}{q}{d} \tt{h}{a}{m} \tt{h}{r}{b}  +2K^{\; \; am}_{[c} K_{d]bm} \\
    & \lera{ \text{$ r \leftrightarrow m $}} \\
    & \tt{ \mathcal{R}}{a}{bcd}  =  \tt{R}{m}{npq} \tt{h}{p}{c } \tt{h}{q}{d} \tt{h}{a}{m} \tt{h}{n}{b}  +2K^{\; \; am}_{[c} K_{d]bm} \\
    & \lera{ \text{$ r \rightarrow n , n \rightarrow p$}} \\  
    &  \tt{ \mathcal{R}}{a}{bcd}  =  \tt{R}{m}{npq} \tt{h}{a}{m}\tt{h}{n}{b}\tt{h}{p}{c } \tt{h}{q}{d}    +2K^{\; \; am}_{[c} K_{d]bm} \\
\end{align*}


And let me put this in a box, as I have to memorize this as the \textcolor{blue}{Make the box. Then start deriving the second Gauss Codazzi relation.}

The second Gauss Codazzi equation is $ \tt{h}{m}{[n} \tt{h}{a}{p]} \tt{h}{b}{q} \tt{k}{c}{r} \nabla_m \tt{K}{}{abc} = \frac{1}{2} \tt{h}{m}{n} \tt{h}{a}{p} \tt{h}{b}{q} \tt{h}{c}{r} \tt{R}{}{mabc}$


To prove it, let's start with the LHS,

\begin{align*}
    LHS &= \tt{h}{m}{[n} \tt{h}{a}{p]} \tt{h}{b}{q} \tt{k}{c}{r} \nabla_m \tt{K}{}{abc} \\
    &= \tt{h}{m}{[n} \tt{h}{a}{p]} \tt{h}{b}{q} \tt{k}{c}{r} \nabla_m \lera{ - \tt{h}{s}{a} \tt{h}{t}{b} \nabla_s h_{tc}}\\
    &= -\tt{h}{m}{[n} \tt{h}{a}{p]} \tt{h}{b}{q} \tt{k}{c}{r} \lera{ \nabla_m \tt{h}{s}{a}} \tt{h}{t}{b} \nabla_s \tt{h}{}{tc} - \tt{h}{m}{[n} \tt{h}{a}{p]} \tt{h}{b}{q} \tt{k}{c}{r} \tt{h}{s}{a} \lera{ \nabla_m\tt{h}{t}{b}} \nabla_s \tt{h}{}{tc} \\
    &- \tt{h}{m}{[n} \tt{h}{a}{p]} \tt{h}{b}{q} \tt{k}{c}{r} \tt{h}{s}{a} \tt{h}{t}{b} \nabla_m \nabla_s \tt{h}{}{tc} \\
    & \equiv -A-B-C \text{  (Naming the three terms.)}\\
\end{align*}


Term A:

\begin{align*}
    A &= \tt{h}{m}{[n} \tt{h}{a}{p]} \tt{h}{b}{q} \tt{k}{c}{r} \lera{ \nabla_m \tt{h}{s}{a}} \tt{h}{t}{b} \nabla_s \tt{h}{}{tc}\\
    &= \tt{h}{m}{[n} \tt{h}{a}{p]} \tt{h}{b}{q} \tt{k}{c}{r} \lera{ \nabla_m \tt{h}{}{as}} \tt{h}{t}{b} \nabla^s \tt{h}{}{tc} \\
    &= \tt{h}{m}{[n} \tt{h}{a}{p]} \tt{h}{t}{q} \tt{k}{c}{r} \lera{ \nabla_m \tt{h}{}{as}} \lera{ \nabla^s \tt{h}{}{tc}} \\
    &= - \tt{K}{}{[np]s} \tt{h}{t}{q} \tt{k}{c}{r} \nabla^s \tt{h}{}{tc} \\
    &= 0 \\
\end{align*}


Term B:


\begin{align*}
    B &=  \tt{h}{m}{[n} \tt{h}{a}{p]} \tt{h}{b}{q} \tt{k}{c}{r} \tt{h}{s}{a} \lera{ \nabla_m\tt{h}{t}{b}} \nabla_s \tt{h}{}{tc} \\
    &= \tt{h}{m}{[n} \tt{h}{s}{p]} \tt{h}{b}{q} \tt{k}{c}{r}  \lera{ \nabla_m\tt{h}{t}{b}} \nabla_s \tt{h}{}{tc} \\
    &= \tt{h}{m}{[n} \tt{h}{s}{p]} \tt{h}{b}{q} \tt{k}{c}{r}  \lera{ \nabla_m\tt{h}{}{bt}} \nabla_s \tt{h}{t}{c} \\
    &= -K_{[n|qt|} \tt{h}{s}{p]} \tt{k}{c}{r} \nabla_s \tt{h}{t}{c} \\
    &= - \tt{k}{u}{t} \tt{K}{}{[n|qu|} \tt{h}{s}{p]} \tt{k}{c}{r}\nabla_s \tt{h}{t}{c} \\
    & \lera{ \text{K is normal in the third index, so introducing k makes sense.}}\\
    &= - \tt{k}{tu}{} \tt{K}{}{[n|qu|} \tt{h}{s}{p]} \tt{k}{c}{r}\nabla_s \tt{h}{}{tc} \\ 
    &= - \tt{k}{t}{u} K^{\;\;\;u}_{[n\;\;\;|q|} \tt{h}{s}{p]} \tt{k}{c}{r} \cov{ \tt{h}{}{tc}}{s} \\
    &= - K^{\;\;\;u}_{[n\;\;\;|q|} \tt{h}{s}{p]} \tt{k}{t}{u} \tt{k}{c}{r} \cov{ \tt{h}{}{tc}}{s} \\
    &= 0 \\
\end{align*}





Term C:

\begin{align*}
   C&=  \tt{h}{m}{[n} \tt{h}{a}{p]} \tt{h}{b}{q} \tt{k}{c}{r} \tt{h}{s}{a} \tt{h}{t}{b} \nabla_m \nabla_s \tt{h}{}{tc} \\
   &= \tt{h}{m}{[n} \tt{h}{s}{p]} \lera{ \nabla_m \nabla_s \tt{h}{}{tc}} \tt{h}{b}{q} \tt{h}{t}{b} \tt{k}{c}{r} \\
   &= \tt{h}{m}{n} \tt{h}{s}{p} \lera{ \nabla_{[m} \nabla_{s]} \tt{h}{}{tc}} \tt{h}{b}{q} \tt{h}{t}{b} \tt{k}{c}{r} \\
   &= \tt{h}{m}{n} \tt{h}{s}{p}\tt{h}{b}{q} \tt{h}{t}{b} \tt{k}{c}{r}\lera{ \nabla_{[m} \nabla_{s]} \tt{h}{}{tc}} \\
   &= \frac{1}{2}  \tt{h}{m}{n} \tt{h}{s}{p}\tt{h}{b}{q} \tt{h}{t}{b} \tt{k}{c}{r} \lera{- \tt{R}{u}{tms} \tt{h}{}{uc}- \tt{R}{u}{cms} \tt{h}{}{tu}} \\
   &= -\frac{1}{2} \lera{\tt{h}{m}{n} \tt{h}{s}{p}\tt{h}{b}{q} \tt{h}{t}{b} \lera{\tt{k}{c}{r} \tt{h}{u}{c}} \tt{R}{}{utms} +\tt{h}{m}{n} \tt{h}{s}{p}\tt{h}{b}{q} \tt{h}{t}{b} \tt{k}{c}{r} \tt{h}{u}{t} \tt{R}{}{ucms}} \\
   &= -\frac{1}{2} \lera{ 0 + \tt{h}{m}{n} \tt{h}{s}{p} \tt{h}{u}{q} \tt{k}{c}{r} \tt{R}{}{ucms} } \\
   &= - \frac{1}{2} \tt{h}{m}{n} \tt{h}{s}{p} \tt{h}{u}{q} \tt{k}{c}{r} \tt{R}{}{msuc} \\
   &= - \frac{1}{2} \tt{h}{m}{n} \tt{h}{a}{p} \tt{h}{b}{q} \tt{k}{c}{r} \tt{R}{}{mabc} \\
\end{align*}

In short, 
\begin{align*}
    \tt{h}{m}{[n} \tt{h}{a}{p]} \tt{h}{b}{q} \tt{k}{c}{r} \nabla_m \tt{K}{}{abc} &= \frac{1}{2} \tt{h}{m}{n} \tt{h}{a}{p} \tt{h}{b}{q} \tt{k}{c}{r} \tt{R}{}{mabc} \\
\end{align*}


\textcolor{blue}{Make another box for this equation.} 

\begin{comment}

\begin{definition}[Your Definition Name]
    Provide a precise mathematical or conceptual definition here. Use proper notation established in the preamble.
    
    For example: A function $f: \R \to \R$ is continuous at $x_0$ if for every $\eps > 0$, there exists $\delta > 0$ such that $|x - x_0| < \delta$ implies $|f(x) - f(x_0)| < \eps$.
\end{definition}




Piecewise definiion

\[
f(x)=
\begin{cases}
x^2, & x \ge 0 \\
-x, & x < 0
\end{cases}
\]







\begin{example}
    Provide a concrete example that illustrates the definition.
\end{example}

\subsection{Theorem 1.1: Your First Result}

\begin{theorem}[Name of Theorem]
    State the theorem clearly here. 
    
    Suppose conditions (i), (ii), and (iii) hold. Then conclusion follows: $A = B$.
\end{theorem}

\begin{proof}
    Outline the proof strategy, then provide key steps.
    
    Step 1: First, we observe that...
    
    Step 2: By applying [technique], we get...
    
    \begin{equation}
        \mathcal{E} = \int_{\Omega} \rho e \, \dd{V}
        \label{eq:energy_def}
    \end{equation}
    
    Step 3: Combining these results yields the conclusion. \qed
\end{proof}

\begin{remark}
    Important observations or extensions of the theorem can go here. This theorem is particularly useful when...
\end{remark}

\section{Key Mathematical Tools}

\subsection{Notation and Operations}

Let's establish the notation we'll use throughout. We denote:
\begin{itemize}
    \item Vectors: $\bv{u} = (u_1, u_2, u_3)$
    \item Matrices: $\mathbf{A} = \{a_{ij}\}$
    \item Derivatives: $\ddt{u} = \frac{\mathrm{d}u}{\mathrm{d}t}$, $\pd{u}{x} = u_{,x}$
\end{itemize}

\subsection{Important Equations}

An important fundamental equation in this field is:

\begin{equation}
    \boxeq{
        \pd{u}{t} + \nabla \cdot \bv{F}(\bv{u}) = 0
    }
    \label{eq:conservation_law}
\end{equation}

This is a conservation law where:
\begin{itemize}
    \item $u$ is the conserved quantity (scalar or vector)
    \item $\bv{F}(\bv{u})$ is the flux function
    \item The equation states that changes in $u$ equal divergence of the flux
\end{itemize}

Equation \eqref{conservation_law} is fundamental and appears repeatedly throughout the course.

\section{Practical Applications}

\subsection{Application 1: [Your Application]}

Consider a specific example where the concepts apply:

\begin{equation}
    \text{Application 1} \quad \rightarrow \quad \text{Result}
    \label{eq:application1}
\end{equation}

The physical significance is that [explain the meaning].

\subsection{Application 2: [Your Application]}

Another important application involves:

\begin{equation}
    \text{Condition} \quad \Rightarrow \quad \text{Consequence}
    \label{eq:application2}
\end{equation}

This demonstrates how [explain the broader context].

\section{Worked Examples}

\begin{example}[Simple Calculation]
    \label{ex:simple_calc}
    
    Given: $\bv{u} = (1, 2, 3)$ and $\bv{v} = (4, 5, 6)$
    
    Find: $\bv{u} \cdot \bv{v}$
    
    \textit{Solution:}
    \begin{align}
        \bv{u} \cdot \bv{v} &= (1)(4) + (2)(5) + (3)(6) \\
        &= 4 + 10 + 18 \\
        &= 32
    \end{align}
\end{example}

\begin{example}[Complex Problem]
    \label{ex:complex_calc}
    
    This more involved example demonstrates integration of multiple concepts. [Provide detailed walkthrough]
\end{example}

\section{Summary and Key Takeaways}

\begin{itemize}
    \item \textbf{Key Point 1:} [Main concept from section 1]
    \item \textbf{Key Point 2:} [Main concept from section 2]
    \item \textbf{Key Point 3:} [Connection to later material]
\end{itemize}

These foundational concepts will be essential for understanding the more advanced material in subsequent chapters.

\section{Exercises}

\begin{enumerate}
    \item \textbf{Basic:} [Simple computational problem]
    \item \textbf{Intermediate:} [Problem requiring multiple concepts]
    \item \textbf{Advanced:} [Challenge problem extending the theory]
    \item \textbf{Proof:} [Problem asking to prove a related result]
\end{enumerate}

\section{Further Reading}

For deeper understanding of the topics covered in this chapter, consult:
\begin{itemize}
    \item [Reference 1] - Good introduction to [topic]
    \item [Reference 2] - Comprehensive treatment of [topic]
    \item [Reference 3] - Advanced material on [topic]
\end{itemize}

%\newpage  % Uncomment to force new page for next chapter



\end{comment}